% Append after the existing supplementary figure sequence; do not insert in S9.
% Asset source: source_data/shunt_ancestry_gain/figures/weak_channel_linearization_check.pdf
\begin{figure}[p]
\centering
\includegraphics[width=\textwidth]{supplementary/figure_shunt_weak_channels.pdf}
\caption{\textbf{Weak voltage-dependent currents provide a local-linearization check near the passive response.}
\textbf{A}, Focal-shunt and baseline-current-matched injection localization
across doses relative to local input conductance. Solid curves use a fixed
Jacobian around each solved weak-channel steady state; dashed curves show the
passive model at the same electrical calibration. The Jacobian $J$ includes
baseline voltage-gate derivatives. A separately solved somatic current
restores baseline somatic voltage within each linearized intervention.
Shunting attenuates descendant gradients; the matched injection enhances
them, although the unsigned localization index is positive for both.
\textbf{B}, Paired shunt-minus-current localization contrasts. Small points
show eight cell means; larger diamonds and bars show cohort means and 95\%
cell-bootstrap intervals. The three doses all have positive contrasts in
eight of eight cells (two-sided Wilcoxon $P=0.0078$). The 64 channel draws
and focal sites are averaged within each cell and are not independent
replicates. Both panels use $R_a=150\,\Omega\,\mathrm{cm}$,
$R_m=1,000\,\Omega\,\mathrm{cm}^2$ and background conductance equal to
baseline leak. Median maximal Na, K, Ca and HCN conductances are 0.06,
0.10, 0.025 and 0.05 times leak; NMDA is 0.15 times excitatory
conductance. This weak-channel ensemble tests the local calculation,
not strong regenerative responses, channel kinetics or a newly solved
nonlinear equilibrium at each dose.}
\label{fig:supp_weak_channel_linearization}
\end{figure}
